\chapter{建立实践共同体}\label{s:community}

你不必为了教编程而去解决社会上所有的弊病，
但如果你想让人们学到东西，你\emph{确实}得参与
课堂之外发生的事。
这对老师和学习者都适用：
许多散养式老师起初是志愿者或兼职者，
得兼顾许多别的责任。
课堂之外发生的事，对他们的成功同样重要，
对学习者的成功也一样，
所以帮助双方的最好办法，就是培育一个教学共同体。

\begin{aside}{芬兰，以及为什么不能照搬}
  芬兰的学校是世界上最成功的之一，
  但正如 Anu Partanen\index{Partanen, Anu}
  \hreffoot{https://www.theatlantic.com/national/archive/2011/12/what-americans-keep-ignoring-about-finlands-school-success/250564/}{所指出的}，
  它们的成功并非孤立取得的。
  其他国家想采用芬兰教学法的尝试注定要失败，
  除非那些国家也能确保孩子（和他们的父母）安全、
  营养充足、
  并在法庭上得到公平对待~\cite{Sahl2015,Wilk2011}。
  考虑到我们对动机对学习之重要性的了解（\chapref{s:motivation}），这并不令人意外：
  如果人们相信这个系统不可预测、不公平或冷漠，每个人都会表现更差。
\end{aside}

思考教学共同体的一个框架是\gref{g:situated-learning}{情境学习}，
它关注\gref{g:legitimate-peripheral-participation}{合法的边缘性参与}
如何让人成为
\gref{g:community-of-practice}{实践共同体}的成员~\cite{Weng2015}。
把这些术语拆开来看，
实践共同体是一群因对某项活动（比如编织或粒子物理）的兴趣而联结在一起的人。
合法的边缘性参与，是指做一些简单的、低风险的任务，
而这些任务被共同体认可为有效的贡献：
织出你的第一条围巾、
在竞选期间装信封、
或为开源软件校对文档。

情境学习关注的是从新来者
到被已有共同体成员接纳为同行的转变过程。
这通常意味着从简单的任务和工具开始，
然后用更复杂的工具做类似的任务，
最后去做和进阶实践者同样的工作。
例如，
学音乐的儿童可能先从在竖笛或尤克里里上弹儿歌开始，
然后加入乐队，用小号或萨克斯吹别的简单曲子，
最后开始探索自己的音乐品味。
支持这一进程的常见方式包括：

\newpage
\begin{description}

\item[解决问题：]
  「我卡住了——我们能一起设计这节课吗？」

\item[请求信息：]
  「邮件列表管理器的密码是什么？」

\item[寻求经验：]
  「有人遇到过有阅读障碍的学习者吗？」

\item[分享资源：]
  「我去年为一个班做了个网站，你们可以拿它当起点。」

\item[协调：]
  「我们能合订 T 恤一起拿折扣吗？」

\item[构建论据：]
  「如果我知道别的集训营是怎么做的，就更容易说服我的老板做出改变。」

\item[记录项目：]
  「这个问题到现在已经出现五次了。让我们一次性把它写清楚。」

\item[梳理知识：]
  「附近街区或城市里还有哪些群体在做类似的事？」

\item[拜访：]
  「我们能来参观你们的课后项目吗？我们需要在自己的城市办一个。」

\end{description}

\hreffoot{https://www.feverbee.com/types-of-community-and-activity-within-the-community/}{大致说来}，
一个实践共同体可以是：

\begin{description}

\item[行动共同体：]
  聚焦于某个共同目标的人，
  比如把某人选上台。

\item[关切共同体：]
  成员因某个共同议题而聚到一起，
  比如应对一种长期疾病。

\item[兴趣共同体：]
  聚焦于对某样东西的共同喜爱，比如西洋双陆棋或编织。

\item[地点共同体：]
  碰巧在同一个地方生活或工作的人。

\end{description}

大多数共同体是这几类的混合，
比如多伦多那些喜欢教技术的人。
一个共同体的焦点也会随时间改变：
例如，
一个帮助抑郁症患者的互助小组（关切共同体）
可以决定筹款来维持一条求助热线（行动共同体）。
而运营这条热线，随后又可以成为这个群体的焦点（兴趣共同体）。

\begin{aside}{先有汤，再有圣歌}
  写宣言很好玩，
  但大多数人加入志愿者共同体是为了帮助别人、也得到帮助，
  而不是为了争论愿景声明的措辞\footnote{
    偏爱后者的人，往往\emph{只}对争论感兴趣{\ldots}}。
  所以你应该聚焦于
  人们能创造出什么、并且能马上被其他共同体成员用上。
  一旦你的组织表明它能做成小事，
  人们就会更有信心，觉得值得帮你做更大的项目。
  那才是该操心定义指导成员价值观的时候。
\end{aside}

\seclbl{先学，再做}{s:community-learn-then-do}

建立共同体的第一步，是决定你是否该建，
还是加入一个已有的组织会更有效。
已经有成千上万的群体在教人们技术技能，
从\hreffoot{http://www.4-h-canada.ca/}{4-H 俱乐部}
和\hreffoot{https://www.frontiercollege.ca/}{扫盲项目}，
到像
\hreffoot{http://www.blackgirlscode.com/}{Black Girls Code}
和\hreffoot{http://bridgeschool.io/}{Bridge} 这样的「带你入门编程」非营利组织。
加入一个已有的群体，会让你在教学上先人一步、
马上有一批同行、
并有机会了解更多运营之道；
但愿，一边做出即时贡献、一边学到这些技能，
比能说自己是某个新事物的创办人或领导者更重要。

不管你是加入已有群体，还是自己办一个，
如果你先读一点关于社区组织工作的背景资料，都会更有效。
\cite{Alin1989,Lake2018}大概是最著名的草根组织工作著作，
而~\cite{Brow2007,Midw2010,Lake2018}是基于数十年实践写成的实用手册。
如果你想读得更深，
\cite{Adam1975}是高地民谣学校（Highlander Folk School）的历史，
它的做法被许多成功的群体效仿；
而~\cite{Spal2014}是一份教成人的指南，
作者在组织工作上有深厚的个人根基；
\hreffoot{https://www.nonprofitready.org/}{NonprofitReady.org}
则提供免费的专业发展培训。

\seclbl{四个步骤}{s:community-four-steps}

每个参与到你的组织里的人
（包括你自己）
都会经历四个阶段：
招募、入职、留存、退出。
刚开始时你不必操心这个循环，
但一旦有超过一小撮的非创办者参与进来，
就值得想一想它了。

第一步是招募志愿者。\index{招募（成员）}
你的市场营销应该在这方面帮到你，办法是让你的组织可被发现，\index{可发现性!组织}
并把它明确的使命和价值传达给
可能想参与的人
（\chapref{s:outreach}）。
分享那些体现你想要的那种帮助的故事，
也分享关于你正在帮助的人的故事，
并说清楚参与的方式有很多。
（我们会在下一节更详细地讨论这一点。）

新成员最好的来源，是你自己的课堂：
「看一个、做一个、教一个」这句口诀，
从志愿者组织\emph{存在}以来就一直对它们很管用。
确保每一次课程或其他接触，
都以告诉人们他们能怎样帮忙、以及他们的帮助会受欢迎来结束。
以这种方式来的人会知道你在做什么，
并且最近才亲身体验过你所提供的东西，
这有助于你的组织避免集体的专家盲区（\chapref{s:memory}）。

\begin{aside}{从小处开始}
  \hreffoot{https://en.wikipedia.org/wiki/Ben\_Franklin\_effect}{本杰明·富兰克林}观察到，
  一个帮过别人忙的人，
  比一个接受过别人帮忙的人，更可能再帮那个人一次忙。
  所以请人帮你做一件小事，
  是让他们愿意帮你做更大事情的好一步。
  教学时一种自然的做法，
  是请人们为你的课程材料提交错字或表述不清处的修正，
  或者建议新的练习或例子。
  如果你的材料是以可维护的方式写的（\secref{s:process-maintainability}），
  这就给了他们一个练习某些有用技能的机会，
  也给了你一个开启对话、进而可能招到一个新成员的机会。
\end{aside}

志愿者生命周期中间的部分，是入职和留存，
我们会在第~\ref{s:community-onboarding}~节和第~\ref{s:community-retention}~节里讲。
最后一步是退出：\index{退出（成员）}
每个人最终都会离开，
而健康的组织会为此做规划。
一些简单的事，就能让离开的人和留下的人都对这次变化感觉良好：

\begin{description}

\item[请人们明确说明他们的离开，]
  好让每个人都知道他们确实已经离开了。

\item[确保他们不为离开或别的事感到难堪或羞愧。]

\item[给他们一个把知识传下去的机会。]
  例如，
  你可以请他们在最后做贡献的几周里指导某个人，
  或者让一位留下的成员访谈他们，
  收集任何值得讲给后人听的故事。

\item[确保他们交还钥匙。]
  某个人离开半年后你才发现，
  只有他一个人知道怎么为一年一度的野餐预订场地，那就很尴尬了。

\item[在他们离开两三个月后再跟进一次，]
  看看他们对你这里什么事管用、什么事不管用还有没有别的想法，
  或者有没有什么建议是他们当时没想到要提、
  或临走时不好意思提的。

\item[感谢他们，]
  既在他们离开时，也在你的群体下次聚会时。

\end{description}

\begin{aside}{缺失的手册}
  关于怎么创办一家公司，已经写了几千本书。
  只有寥寥几本讲怎么优雅地结束一家公司、或离开一家公司，
  尽管每一个开始都有一个结束。
  如果你哪天写了一本，
  请告诉我。
\end{aside}

\seclbl{入职}{s:community-onboarding}

决定加入一个群体之后，
人们需要尽快上手，
而~\cite{Shol2019}总结了关于这件事我们知道些什么。\index{入职（成员）}
第一条规则是有并执行一份行为准则（\secref{s:classroom-coc}），\index{行为准则}
并找到一个愿意接收和审查不当行为报告的独立第三方。
组织外部的人会有组织成员可能缺乏的客观性，
也能保护那些担心报复或名誉受损、
因而犹豫是否就项目领导者提出问题的人。
项目领导者还应该公开执行决定，
好让共同体认识到这份准则是有分量的。

接下来最重要的事，是友善待人。
正如 Fogel 所说~\cite{Foge2005}，
「如果一个项目没有给人留下良好的第一印象，
新来者可能要等很久才会给它第二次机会。」
其他作者也在实证上证实了友善礼貌的社交环境在开放项目中的重要性~\cite{Sing2012,Stei2013,Stei2018}：

\begin{description}

\item[发布一条欢迎信息，]
  发在项目社交媒体页面、Slack 频道、论坛或邮件列表上。
  项目可以考虑维护一个专门的「欢迎」频道或列表，
  由项目负责人或社区经理写一条简短帖子，请新来者自我介绍。

\item[帮人们找到做首次贡献的方式，]
  比如把某些需要加工的课程或工作坊标注为「适合新来者」，
  并请老成员不要去修它们，
  以保证新来者有合适的起点。

\item[把新来者引向和他们相似的项目成员，]
  以表明他们属于这里。

\item[给新来者指出项目的关键资源，]
  比如贡献指南。

\item[指定一两名成员作为联系人，]
  负责接待每位新来者。
  这么做能让新来者不那么不愿提问。

\end{description}

有助于所有人（不只是新来者）的第三条规则，
是让知识可被找到、并保持更新。
新来者就像探险者，必须在陌生的景色里给自己定位~\cite{Dage2010}。
信息如果分散在各处，通常会让新来者感到迷失和无所适从。
考虑到维护信息的地方有很多种可能
（例如 wiki、版本控制里的文件、共享文档、旧的推文或 Slack 消息、以及邮件存档），
重要的是把关于某个特定主题的信息整合在一个地方，
这样新来者就不必穿梭于多个数据源去找他们需要的东西。
把信息组织好，能让新来者更有信心、更清楚方向~\cite{Stei2016}。

最后，
肯定新来者的首次贡献，
并弄清他们长期来看可以在哪里、怎样帮上忙。
一旦他们把自己的首次贡献过了线，
你和他们就都可能更清楚他们能提供什么、以及项目能怎样帮到他们。
帮新来者找到下一个他们可能想上手的问题，
或给他们指出下一件他们可能会喜欢读的东西。
尤其是，
鼓励他们去帮下一批新来者，
既是认可他们所学东西的好办法，
也是把它传下去的有效方式。

\seclbl{留存}{s:community-retention}
\index{留存（成员）}

\begin{quote}

  如果你的人做这件事时一点乐子都没有，那就有大问题了。\\
  —— Saul Alinsky\index{Alinsky, Saul}

\end{quote}

共同体成员不该指望在为你组织工作的每一刻都享受，
但如果他们一点都享受不到，
他们就不会留下来。
享受不一定意味着每年办一场派对：
人们可能享受做饭、
带教、
或只是安静地坐在别人旁边工作。
每个组织都应该做一些事，来确保
人们能从工作中得到他们看重的东西：

\begin{description}

\item[去问人们想要什么，而不是猜。]
  正如你不是你的学习者（\secref{s:process-personas}）一样，
  你和组织里其他成员大概也不一样。
  问人们想做什么、
  愿意做什么（这可能不是一回事）、
  以及他们的时间有什么约束。
  他们可能会说「什么都行」，
  但哪怕只是简短聊一聊，大概也能发现
  他们喜欢和人打交道，却不太愿意管小组的财务，
  或者反过来。

\item[提供多种贡献方式。]
  让人出力的方式越多，
  能出力的人就越多。
  一个不喜欢站在听众面前的人，
  也许能维护你组织的网站、
  处理账务、
  或校对课程。

\item[认可贡献。]
  人人都喜欢被欣赏，
  所以共同体应该公开和私下都认可
  成员的贡献，
  比如在演示里提到他们、
  把名字放到网站上，等等。
  某人为你的项目付出的每一个小时，
  都可能是从他的个人生活或本职工作中挤出来的一小时；
  要认识到这一点，
  并明确表示：虽然多出力欢迎，
  但你并不指望他们做出无法持续的牺牲。

\item[留出空间。]
  你以为自己在帮忙，
  但插手每一个决定，会剥夺人们的自主感，
  这反过来又降低他们的动机（\secref{s:motivation}）。
  尤其是，
  如果你总是第一个回复邮件或聊天消息，
  人们就少了成长为成员的机会，
  也少了建立横向协作的机会。
  结果，
  共同体就会继续围绕一两个人转，
  而不是变成一个联系紧密、别的人也乐于参与的网络。

\end{description}

奖励参与的另一种方式，是提供培训。
组织需要预算、拨款申请和纠纷解决。
大多数人从没被教过怎么做这些，就像他们从没被教过怎么教书一样，
所以获得可迁移技能的机会，
是人们加入并留下来的一个强有力理由。
如果你要这么做，
除非这正是你的专长，否则不要试图自己提供培训。
许多公民和社区团体有这类项目，
你大概能和其中一个达成协议。

最后，
虽然志愿者能做很多事，
但像系统管理和会计这样的任务，最终需要带薪员工。
当你到了这一步，
要么一分钱不付，要么付一份像样的工资。
如果你一分钱不付，
他们真正的回报就是行善的满足感；
另一方面，如果你只付象征性的钱，
你就把那份满足感拿走了，却没有给他们谋生的满足感。

\seclbl{治理}{s:community-governance}
\index{治理}

每个组织都有一个权力结构：
唯一的问题是
它是正式且可问责的，还是非正式因而不可问责的~\cite{Free1972}。
后者对多则半打人、且彼此都认识的小组其实相当管用。
超过这个规模，
你就需要规则来写清楚
谁有权做哪些决定、
以及如何达成共识（\secref{s:meetings-marthas-rules}）。

我更偏好的治理模式是\gref{g:commons}{共有资源}，
即由某个社群按照他们自己演化并采纳的规则共同管理的东西~\cite{Ostr2015}。
正如~\cite{Boll2014}强调的，
这个定义的三部分缺一不可：
共有资源不只是共用的牧场或一组软件库，
它还包括共享它的共同体
以及他们用来共享的规则。

营利性公司和注册的非营利组织是更流行的模式；
具体机制因司法管辖而异，
所以做选择之前你应当先咨询\footnote{
  这正是和地方政府或其他志同道合的组织保持联系能派上用场的时刻之一。}。
这两类组织都把最终权力交给自己的董事会。
大致说来，这要么是一个\gref{g:service-board}{服务型委员会}，
其成员还在组织里承担别的角色；
要么是一个\gref{g:governance-board}{治理委员会}，其主要职责是聘用、监督，
必要时解雇负责人。
董事会成员可以由共同体选举产生，也可以被任命；
无论哪种情况，
重要的是把胜任放在热情之上
（后者对普通成员更重要），
并设法招募具备特定技能的人，比如会计、市场营销等。

\begin{aside}{选择民主}
  时机到了的时候，
  把你的组织变成民主制：
  或早或晚（通常偏早），
  每一个被任命的董事会都会变成互相抬轿子的圈子。
  把权力交给成员会很混乱，
  但这是迄今发明出来的、能确保组织持续满足人们实际需求的唯一办法。
\end{aside}

\seclbl{照顾好你自己}{s:community-care}

倦怠是任何社区活动里的慢性风险~\cite{Pign2016}，\index{倦怠}
所以要学会说「不」多于说「是」。
如果你不照顾好自己，
就照顾不了你的共同体。

\begin{aside}{「不」用光了}
  1990 年代的研究似乎表明，我们施展意志力的能力是有限的：
  如果我们饿着肚子还得忍住不吃盒里最后那个甜甜圈，
  就更不可能去叠衣服，反之亦然。
  这一现象叫\gref{g:ego-depletion}{自我损耗}，
  尽管近期研究没能重复出那些早期结果~\cite{Hagg2016}，
  但在太累而说不出「不」的时候说「是」，
  是许多组织者会掉进去的陷阱。
\end{aside}

让你的「不」站得住的一种办法，
是写一份「不做清单」，列出那些值得做、
但你\emph{不打算}做的事。
写这本书时，
我的清单上有四本书、
两个软件项目、
重新设计我的个人网站、
以及学吹锡口笛。

最后，
时不时提醒自己：
每个组织最终都需要新想法和新领导。
到那时候，
培养你的继任者，然后尽量优雅地离开。
他们无疑会做一些你不会做的事，
但生活中很少有什么事，能比看着你参与建成的东西有了自己的生命更让人满足。
为此庆祝吧——你不会有任何困难去找别的事让自己忙起来。

\seclbl{练习}{s:community-exercises}

其中几个练习取自~\cite{Brow2007}。

\exercise{哪种共同体？}{个人}{15}

重读四种共同体的描述，
判断你的群体是或想成为哪一种（或几种）。

\exercise{你可能会遇到的人}{小组}{30}

作为组织者，
你的一部分工作有时是帮人们找到一种方式，让他们即使有自身特点也能做出贡献。
以小组为单位，
从下面这些人里挑三个，
讨论你会怎样帮他们成为你组织中更好的贡献者。

\begin{description}

\item[Anna]
  对每个主题的了解都比其他所有人加起来还多——至少，
  她自己这么认为。
  不管你说什么，
  她都会纠正你；
  不管你懂什么，她都懂得更多。

\item[Catherine]
  对自己的能力太没信心，
  以至于不管多小的决定，
  都要先问过别人，才敢做。

\item[Frank]
  喜欢知道别人不知道的事。
  他能创造奇迹，
  但别人问他是怎么做到的，
  他就咧嘴一笑说：
  「哦，我相信你能自己想出来。」

\item[Hediyeh]
  很安静。
  她在会上从不发言，
  即便知道别人错了。
  她可能会在邮件列表上出力，
  但她对批评非常敏感，
  总是退让，而不是为自己的观点辩护。

\item[Ken]
  利用了大多数人宁愿分担他那份工作、
  也不愿抱怨他的事实。
  让人恼火的是，终于有人当面质问他时，他说得那么像回事。
  「各方面都有错误，」他说，
  或者，「唉，我觉得你在鸡蛋里挑骨头。」

\item[Melissa]
  本意是好的，
  但不知怎的，总有事冒出来，
  她的任务永远是在最后一刻之前都完不成。
  当然，
  这意味着所有依赖她的人，
  都得等到最后一刻\emph{之后}才能做自己的事{\ldots}

\item[Raj]
  很粗鲁。
  「我就是这么说话的，」他说。
  「你要是受不了，就另找团队吧。」
  他最爱说的一句话是「那很蠢」，
  而且每隔一句话就带一个脏字。

\end{description}

\exercise{价值观}{小组}{45}

先自己回答这些问题，
再和别人比较答案。

\begin{enumerate}

\item
  你的组织表达了哪些价值观？

\item
  这些是你希望组织表达的价值观吗？

\item
  如果不是，你希望它表达什么价值观？

\item
  有哪些具体行为体现了这些价值观？

\item
  哪些行为会体现与这些价值观相反的东西？
\end{enumerate}

\exercise{会议程序}{小组}{30}

先自己回答这些问题，
再和别人比较答案。

\begin{enumerate}

\item
  你们的会议是怎么开的？

\item
  这是你希望会议开成的方式吗？

\item
  开会的规则是明确写出来的，还是只是默认的？

\item
  这些是你想要的规则吗？

\item
  谁有资格投票或做决定？

\item
  这是你希望被赋予决策权的人吗？

\item
  你们用多数决、共识决，还是别的什么？

\item
  这是你希望做决定的方式吗？

\item
  会上的人怎么知道一个决定已经做出？

\item
  没参会的人怎么知道做了哪些决定？

\item
  这对你的群体管用吗？

\end{enumerate}

\exercise{规模}{小组}{20}

先自己回答这些问题，
再和别人比较答案。

\begin{enumerate}

\item
  你的群体有多大？

\item
  这是你希望组织的规模吗？

\item
  如果不是，你希望它有多大？

\item
  你们对成员规模有任何限制吗？

\item
  设定这样一条限制对你们有好处吗？

\end{enumerate}

\exercise{成为成员}{小组}{45}

先自己回答这些问题，
再和别人比较答案。

\begin{enumerate}

\item
  人们怎么加入你的群体？

\item
  这个过程效果如何？

\item
  要交会费吗？

\item
  加入时要求人们同意某些行为规则吗？

\item
  这些是你想要的行为规则吗？

\item
  新来者怎么知道需要做什么？

\item
  这个过程效果如何？
  
\end{enumerate}

\exercise{人员配置}{小组}{30}

先自己回答这些问题，
再和别人比较答案。

\begin{enumerate}

\item
  你的组织有带薪员工吗，
  还是所有人都是志愿者？

\item
  你应该有带薪员工吗？

\item
  你希望/需要更多还是更少的员工？

\item
  员工们都做什么？

\item
  这些是你需要员工来承担的主要角色和职能吗？

\item
  谁监督你的员工？

\item
  这是你希望你的群体有的监督流程吗？

\item
  你的员工拿多少报酬？

\item
  这是完成所需工作的合适薪水吗？

\end{enumerate}

\exercise{钱}{小组}{30}

先自己回答这些问题，
再和别人比较答案。

\begin{enumerate}

\item
  谁为什么付钱？

\item
  这是你希望付钱的人吗？

\item
  你的钱从哪里来？

\item
  这是你希望获得钱的方式吗？

\item
  如果不是，你有用别的办法搞到钱的计划吗？

\item
  如果有，是什么？

\item
  谁在跟进，以确保这件事发生？

\item
  你有多少钱？

\item
  你需要多少？

\item
  你的钱大部分花在哪里？

\item
  这是你希望花钱的方式吗？

\end{enumerate}

\exercise{借鉴想法}{全班}{15}

我关于如何建立共同体的许多想法，
都受到我在开源软件开发中经历的塑造。
\cite{Foge2005}（可在网上\hreffoot{http://producingoss.com/}{免费获取}）
是一份很好的指南，讲那些共同体里什么管用、什么不管用；
\hreffoot{https://opensource.guide/}{开源指南网站}
也有大量有用信息。
从后者中选一节，
比如「为你的项目找到用户」
或「领导与治理」，
在小组里做两分钟展示，讲其中一个你觉得有用、
或你强烈反对的想法。

\exercise{你是谁？}{小组}{20}

美国国家海洋与大气管理局（NOAA）
有一份简短、有用又风趣的指南，讲
\hreffoot{https://coast.noaa.gov/ddb/}{如何应对破坏性行为}。
它把这些行为分成「话多」「优柔寡断」「害羞」之类的标签，
并概述了应对每一种的策略。
以 3—6 人为一组，
读这份指南，判断哪种描述最符合你。
你觉得为应对像你这样的人所描述的策略有效吗？
别的策略是否同样有效或更有效？

\exercise{一起创作课程}{小组}{30}

\hreffoot{http://carpentries.org}{The Carpentries} 成功的关键之一，
是它们强调协作地构建和维护课程~\cite{Wils2016,Deve2018}。
以 3—4 人为一组：

\begin{enumerate}

\item
  挑一节你们都上过的短课。

\item
  认真审阅，整理出一份统一的改进建议清单。

\item
  把这些建议提供给课程的作者。

\end{enumerate}

\exercise{你「脆」了吗？}{个人}{10}

\hreffoot{https://mailchi.mp/d54702d0a790/take-my-horse-to-the-sand-hill-road}{Johnathan Nightingale 写道}：

\begin{quote}
  我在 Mozilla 工作时，
  我们用「脆」（crispy）这个词来指倦怠临界前的状态。
  变脆的人不好相处。
  他们说话简短生硬。
  他们巴不得找场能赢的架来打。
  他们动不动就哭。
  {\ldots}我们会在同事身上认出这种「脆」，并互相照应，
  [但]我们见得太多、以致围绕它形成了一整套文化流程，这本身是件难堪的事。
\end{quote}

\noindent
对下面每个问题回答「是」或「否」。
你离倦怠有多近？

\begin{itemize}
\item 你在工作中变得愤世嫉俗或爱挑剔了吗？
\item 你是不是得强迫自己去上班，或者很难开始工作？
\item 你是否变得易怒或对同事不耐烦？
\item 你是不是觉得很难集中注意力？
\item 你是否无法从成就中获得满足感？
\item 你是否在用食物、药物或酒精让自己感觉好些、或干脆不要有感觉？
\end{itemize}
